\documentclass[10pt,a4paper]{amsart}
\usepackage[margin=19mm]{geometry}
\usepackage{amsmath,amssymb,mathtools}
\usepackage[scr=rsfs,cal=euler]{mathalfa}
\usepackage{xfrac}
\usepackage[unicode,hidelinks,bookmarks=false]{hyperref}
\newcommand{\ie}{i.e.\ }
\newcommand{\cf}{cf.\ }
\newcommand{\resp}{respectively\ }
\newcommand{\Subsection}{\subsection}
\providecommand{\nobreakdash}{\nobreak\hbox{-}}
\setlength{\emergencystretch}{2em}
\multlinegap0pt
\usepackage{amssymb}
\usepackage{mathtools}
\usepackage[shortlabels]{enumitem}
\usepackage{tikz}
\newtheorem{rem}{Remark}
\newtheorem{lem}{Lemma}
\newtheorem{thm}{Theorem}
\usepackage{cleveref}
\crefname{rem}{Remark}{Remarks}
\crefname{lem}{Lemma}{Lemmas}
\crefname{thm}{Theorem}{Theorems}
\newcommand{\dom}[1]{\mathbf{dom}\left(#1\right)}
\newcommand{\ran}[1]{\mathbf{ran}\left(#1\right)}
\title{Geometric Inverse Semigroup Theory: a note on the Milnor-Schwarz Lemma for inverse monoids}
\author{Giorgio Mangioni, Francesco Tesolin}
\date{Source: arXiv:2603.29524v2 (CC BY 4.0)}
\begin{document}
\maketitle

\noindent\emph{Source and licence.} Geometric Inverse Semigroup Theory: a note on the Milnor-Schwarz Lemma for inverse monoids. Version arXiv:2603.29524v2 (CC BY 4.0), distributed by arXiv under the Creative Commons Attribution 4.0 International licence, http://creativecommons.org/licenses/by/4.0/, as recorded in the arXiv licence metadata for this exact version and shown on the abstract page https://arxiv.org/abs/2603.29524v2. Full licence notice: \texttt{licenses/arxiv-2603.29524-CC-BY-4.0.txt}.\par\smallskip

\noindent\textit{Editorial scope.} Counted unit: the article's thm thm:MilnorSchwarz (source0.tex lines 377--384). The article's complete original proof of it is delivered with it (source0.tex lines 386--440, one \texttt{proof} environment of 55 lines). No mathematics is written, repaired or re-derived: the statement and the proof are the article's own lines, in the article's own notation, and no retained line is substituted or re-flowed.  Retained uncounted as the source-local context the proof needs: lines 200--203; lines 280--283; lines 321--326.  Nothing was omitted from any retained span, so the package carries no omission record.  No retained line contains a citation, so the wrapper carries no bibliography and no bibliography entry.  No retained line contains a figure environment, an image-inclusion command or drawing code, so no image is delivered and the package declares no asset dependency.  Editorial formatting only, in addition: an AMS \texttt{amsart} preamble that restates the article's own declarations and macros used by the retained lines, namely the environments \texttt{rem}, \texttt{lem}, \texttt{thm} on separate counters, together with the standard packages those lines need.  The retained environments print in this document as Remark 1, Lemma 1, Theorem 1 in order of appearance, whereas the article prints the same items with the numbering of its own full document.  The complete native source of the article as distributed in the arXiv e-print for this exact version is bundled unchanged as \texttt{sources/197518/source0.tex}.
\begin{rem}\label{rem:order_under_products}
	Given $s,t\in S$, we have that $\dom{st}\le \dom{t}$, since
	$$\dom{st}\dom{t}=t^{-1}s^{-1}stt^{-1}t=t^{-1}s^{-1}st=\dom{st}.$$
\end{rem}

\begin{lem} \label{lem:s_induces_isom}
    Let $S$ act on $(X,E(S),p)$ a presheaf of metric spaces. Each $s \in S$ induces an isometry
    $\theta_s \colon X \cdot ss^{-1} \to X \cdot s^{-1}s$ by acting on the right by $s$.
\end{lem}

\begin{lem}\label{lem:cayley_dist_without_idempotents}
    Let $\mathcal M$ be a quasi-generating set for $S$. Then $d_{\mathcal M\cup E(S)}=d_{\mathcal{M}}$, where $$d_{\mathcal{M}}(s,t)=\begin{cases}
        \min\{k\mid \exists m_1,\ldots, m_k\in \mathcal{M} \mbox{ s.t. }s= m_1\ldots m_kt\}&\mbox{if }(s,t)\in \mathcal{L};\\
        \infty &\mbox{if }(s,t)\not\in \mathcal{L}.
    \end{cases}$$
\end{lem}

\begin{thm}\label{thm:MilnorSchwarz}
	Let $S$ be an inverse monoid and $(X, E(S),p)$ be a presheaf of geodesic metric spaces. 
	If $S$ acts on $X$ properly and coboundedly then:
    \begin{enumerate}
        \item There exists a finite, quasi-generating set $\mathcal{M}\subseteq S$.
        \item\label{item:MS:qi} There exists $x_1\in X_1$ such that the orbit map $(S,d_{\mathcal{M}})\to X$ sending $s\in S$ to $x_1\cdot s$ is an order-preserving quasi-isometry.
    \end{enumerate}
\end{thm}

\begin{proof}
	By coboundedness there exist an 
    $x_1 \in X_1$ and $T\ge 0$ such that $B(x_1,T) \cdot S = X$. Let $$G \coloneq \{ s \in S \mid d(x_1 \cdot s, x_1\cdot \dom{s})\le 2T+1\},$$ we work to show that $G$ is a generating set of $S$. For every $s \in S$, $X_{\dom{s}}$ is a geodesic metric space, so there exists a geodesic
	$\gamma$ from $x_1 \cdot \dom{s}$ to $x_1 \cdot s$ in  $X_{\dom{s}}$.
    Take points along the geodesic at most a distance $1$ apart: concretely,
    let $$x_1 \cdot \dom{s} = p_1, \,p_2,\, \cdots,\, p_k = x_1 \cdot s$$
    be a series of points
    such that $d(p_i,p_{i+1}) \leq 1$ and $k \leq d(x_1 \cdot \dom{s}, x_1 \cdot s) + 2$.
    Since each $p_{i}$ is in $B(x_1,T) \cdot S$ there exists elements $s_{i}$ in $S$ 
	such that $d(p_{i}, x_1 \cdot s_{i}) \leq T$; for $i=k$ we can choose $s_k=s$. This implies that $x_1 \cdot s_i \in X_{\dom{s}}$ and hence
	$\dom{s_{i}} = \dom{s}$ for all $i$.
	Define $u_1 = s_1$ and $u_{i} = s_{i} s_{i-1}^{-1}$ for $2 \leq i \leq k$; notice that $u_k=ss_{k-1}^{-1}$. Then,
	\[
		s = s\dom{s}^{k-1}=s\dom{s_{k-1}}\ldots\dom{s_1}=(s s_{k-1}^{-1})(s_{k-1}s_{k-2}^{-1})\ldots (s_2s_1^{-1})s_1 = u_{k} \ldots u_{1}.
	\]
    Now, by assumption
    $d(x_1 \cdot s_1, x_1 \cdot \dom{s_1}) \leq T$, so $s_1\in G$. Moreover, for all $2 \leq i \leq k$,
	\[d(x_1 \cdot u_i, x_1 \cdot \dom{u_i}) = d(x_1 \cdot s_is_{i-1}^{-1}, x_1 \cdot s_{i-1}s_i^{-1}s_is_{i-1}^{-1})\]\[ =d(x_1 \cdot s_is_{i-1}^{-1}, x_1 \cdot s_{i-1}s_{i-1}^{-1}) \leq d(x_1\cdot  s_{i}, x_1 \cdot s_{i-1}) \leq 2T + 1,\]
	where we used that $s_i^{-1}s_i=\dom{s_i}=\dom{s_{i-1}}=s_{i-1}^{-1}s_{i-1}$. 
    Hence,  all $u_i$ are in $G$, so $G$ is a generating set.
    By properness $G \subseteq \mathcal{C}E(S)$ where $\mathcal{C}= \mathcal{C}(x_1, 2T+1)$, so
    $\mathcal{M} = \mathcal{C}$ is a finite quasi-generating set of $S$.

    We now work to prove (2). The orbit map $S\to x_1\cdot S$ is $T$-coarsely surjective by coboundedness, so we focus on showing it is a quasi-isometric embedding. Note that $$d_{\mathcal{M}}(s,t) < \infty \iff \dom{s}=\dom{t} \iff
    d(x_1 \cdot s, x_1 \cdot t) < \infty.$$ Hence, if $d_{\mathcal{M}}(s,t)<\infty$ then 
    \[d(x_1 \cdot s, x_1 \cdot t) =d(x_1 \cdot s\dom{t}, x_1 \cdot t\dom{s})= d(x_1 \cdot st^{-1}, x_1 \cdot \dom{st^{-1}}),\]
    where we used \Cref{lem:s_induces_isom} for the last equality. The same argument applies to the Cayley metric $d_\mathcal{M}$; hence it is enough to show that, for every $s\in S$, $d(x_1\cdot \dom{s}, x_1\cdot s)$ is linearly bounded above and below in terms of $d_{\mathcal{M}}(\dom{s}, s)$. 
    
   Write
    $s = u_{k} \cdots u_1$ as before. For each $i$ denote $v_{i} = u_{i} \cdots u_1 \dom{s}$, with $v_0 = \dom{s}$. By \Cref{rem:order_under_products}
\[
	\dom{s} = \dom{v_{k}} 
	\leq \dom{v_{k-1}} \leq \dom{v_{k-2}} \cdots\leq \dom{v_0} = \dom{s}
.\]
Hence, $\dom{v_i} = \dom{s}$ for all $i$, so each $u_i$ labels a (possibly loop) edge in the Schützenberger graph of $s$ connecting $v_{i-1}$ to $v_i$. Therefore, by construction of the $u_i$,
$$d_{\mathcal{M}}( \dom{s}, s) \leq k \leq d(x_1 \cdot \dom{s}, x_1 \cdot s) + 2.$$
For the other inequality, let $n=d_{\mathcal M}(s,\dom{s})$, and write $s=t_n\ldots t_1\dom{s}$, where each $t_i\in \mathcal M$ (this can be done by \Cref{lem:cayley_dist_without_idempotents}). Set $v_i=t_i\ldots t_1 \dom{s}$, with $v_0=\dom{s}$. Note that, for all $i$,
\[\dom{v_{i}} = \dom{v_{i+1}} = \dom{t_{i+1}v_{i}} = v_{i}^{-1}\dom{t_{i+1}}v_{i};\] hence
\[\ran{v_{i}}=v_i \dom{v_i}v_i^{-1}=\ran{v_i}\dom{t_{i+1}}\ran{v_i}= \ran{v_i}\dom{t_{i+1}}\leq \dom{t_{i+1}},\]
where we used that idempotents commute. Therefore, combining \Cref{lem:s_induces_isom} and the fact that the action is by $1$-Lipschitz maps, we get
 \[
	 d(x_1 \cdot v_{i+1}, x_1 \cdot v_{i}) =d(x_1 \cdot t_{i+1}v_i, x_1 \cdot v_{i}) = d(x_1 \cdot t_{i+1} \ran{v_{i}}
	 , x_1 \cdot \ran{v_{i}})
	 \leq d(x_1 \cdot t_{i+1}, x_1 \cdot \dom{t_{i+1}})
.\]
Hence
\[
	d(x_1 \cdot \dom{s}, x_1 \cdot s)
	\leq \sum_{i=1}^{n} d(x_1 \cdot v_{i-1}, x_1 \cdot v_{i})\leq n \max_{i=1,\ldots, n} \{d(x_1 \cdot t_{i}, x_1 \cdot \dom{t_{i}}) \}.\]
Since the $t_i$'s belong to the finite set $\mathcal M$, the quantity $d(x_1 \cdot t_{i}, x_1 \cdot \dom{t_{i}})$ is uniformly bounded, thus concluding the proof.
% Let $t_i=m_i e_i$, where $m_i\in \mathcal M$ and $e_i\in E(S)$. Then $\dom{t_{i}}=\dom{m_i}e_i$, so
% \[
% d(x_1\cdot m_i, x_1\cdot \dom{m_i})\ge d(x_1\cdot t_i, x_1\cdot \dom{m_i}e_i)=d(x_1\cdot t_i, x_1\cdot \dom{t_i}).\]
% Notice that the left-hand-side of the inequality is uniformly bounded since $\mathcal M$ is finite, and we are done.
\end{proof}

\end{document}
